\section{Boundaries and outlook}
\label{sec:outlook}

The parity theorem settles whether literal reverse-complement symmetry can be
achieved at the ordinary necklace bound for every finite alphabet with a
fixed-point-free complement.  Several nearby questions remain distinct.

\begin{enumerate}
  \item The exact premium $\Delta_q(k)$ at odd orders is not determined.
  The proof gives the uniform lower bound $q$ by counting disjoint
  obstructions, one for each complementary alphabet pair. This is a lower
  bound, not an equality claim or a packing--covering theorem.

  \item The descent does not identify a unique terminal or a shortest path.
  The direct selector chooses one terminal after fixing weights and a total
  alphabet order, but different choices may produce different MDSs.

  \item Minimum cardinality implies that the residual graph is acyclic, but
  it does not bound its longest directed path.  Quantitative window guarantees
  require a separate analysis.

  \item The fixed-point-free hypothesis is structural.  If a complement has
  fixed letters, lower reverse-complement fixed points and odd-order PCR
  arguments change.  The exceptional order $k=1$ also behaves differently.

  \item Corollary~\ref{cor:complexity} is a word-membership theorem, not a
  claim that explicitly listing an exponentially large MDS takes polynomial
  total time.

  \item No biological conservation, density, runtime, memory, clinical, or
  application-performance conclusion follows from the graph theorem alone.
\end{enumerate}

The self-dual ideal description may nevertheless be useful independently of
the present proof: it converts a global symmetry condition into an exact
local membership rule while retaining the ordinary optimum.  Understanding
which choices of weights and midpoint ties also optimize residual-path
statistics is a natural next mathematical problem.

\section*{Use of generative AI}
\addcontentsline{toc}{section}{Use of generative AI}

Generative-AI assistants contributed substantively to the mathematical
development, including candidate proof arguments, symbolic identities,
construction and verification code, adversarial checks, and manuscript
drafting. The author selected the problem, directed the research process,
and is responsible for the article's claims and final text. The two exact
implementations are distinct, but the AI-assisted review process does not
constitute independent specialist peer review. No end-to-end formal
proof-assistant certification is claimed.
